\section{Elección de constantes para el corte proporcional}
Se admiten las estimaciones deterministas citadas en el registro de fuentes, con sus
hipótesis uniformes. El objetivo aquí es la compatibilidad de los parámetros.
\begin{enumerate}
\item Se fijan la banda histórica de C, su $\delta$ y la tasa de salidas con su coste.
Se reservan $M_0,\gamma_T$ de la parte 01. Para los lemas locales se toma una clase
de referencia con esos márgenes y una cota superior fija de masa admisible.
Se usa la envolvente de \emph{Grados ponderados}, parte 05, en el alcance
$v\le 8\log n$ (logaritmos naturales). Sus constantes son las de la clase de referencia.
Las caras sólo activan los $v$ sitios en el mismo exterior de torres. Los radios locales
$\kappa/v,\kappa'/d_i$ dependen de los márgenes del exterior, según sus partes 02--04.
No se deduce uniformidad sobre máscaras arbitrarias a partir del cálculo uniforme del hipergrafo.
\item Sea $C_1$ una base común que domine las de la suma puntual y en media, antes
de introducir el coste $\delta^{-2v}$ de prescripción. Elegimos $0<a<1$ con
$a\log(\max(2,C_1))<1/4$ y después $\rho\ge2$ con $a\log\rho>3$.
Los dos cortes $\lfloor a\log j\rfloor$ y $\lceil4\log n/\log\rho\rceil$ caben en
$8\log n$ para $n$ grande. Así el alcance empleado para obtener $C_1$ no se recalcula.
\item Con $\eta$ el umbral común de las dos normas de activación para ese radio,
escogemos $0<\alpha<\alpha_{\rm hist}$ cumpliendo la parte 01 y además
$M_0\alpha<\min(\epsilon_{\mathcal D},\eta/(2\rho^2))$ para la clase ya fijada.
En el extremo corto, $\rho^2(\epsilon_A+1/j)\le\eta$ para $j$ grande.
Si una entrada se expresa en carga de $D_A$, se impone asimismo el umbral
correspondiente con $4C\alpha$. Son constantes fijadas antes que $\alpha$.
Escogemos $\delta_*\in(0,1)$ con $C_1\delta^{-2}\delta_*^{1/3}<1$.
\item Se toma $0<b<\min(\delta/2,\delta_*)$ y se disminuye una sola vez para que
\[
 \frac{\alpha}{4}\log\frac{2b}{\alpha\delta}
        +\frac{C_{\rm pres}\alpha^2}{16}<-(\beta+3),\qquad
 \frac{2b}{\alpha\delta}<1 .
\]
Es posible con $\alpha,\delta,\beta,C_{\rm pres}$ ya fijos. Se usa la estimación de
longitudes de \emph{Llegada y completabilidad}, parte 02. Para $j=\lfloor bn\rfloor$,
$\lceil\sqrt n\rceil\le j$ y $v^2/j=o(1)$ hasta $8\log n$. Todas las elecciones
preceden al límite $n\to\infty$. No se afirma uniformidad de la tasa en $b\downarrow0$.
\end{enumerate}
Bajo estas premisas, los dos radios dan los respectivos restos
$O_\rho(j^2\rho^{-\lfloor a\log j\rfloor})=O_\rho(j^{-1})$ y $O_\rho(n^{-2})$.
La suma de grados de la parte 07 de \emph{Grados ponderados} sigue siendo geométrica
con el factor $\delta^{-2v}$ bajo $\nu$. El suceso corto lleva indicador, no se divide
por su probabilidad. Las estimaciones locales se aplican en ese extremo.
La acotación previa de cada base en una clase fija es esencial. Si tal base
se conociese sólo tras elegir $\alpha$, este orden de elección no la suministraría.
